\chapter{Ideale und Hauptfaktoren}

\section{Idealbildung}

\begin{definition}[Links-, Rechts- und zweiseitige Ideale]
Eine nichtleere Teilmenge $I\subseteq S$ heißt \emph{Linksideal},
wenn $SI\subseteq I$, und \emph{Rechtsideal}, wenn $IS\subseteq I$ gilt.
Ein \emph{Ideal} ist zugleich Links- und Rechtsideal.
Für $X\subseteq S$ heißen $S^1X$, $XS^1$ und $S^1XS^1$
die von $X$ erzeugten Links-, Rechts- und zweiseitigen Idealabschlüsse.
Für $X=\{a\}$ schreibt man $S^1a$, $aS^1$, $S^1aS^1$ und spricht
von \emph{Hauptidealen}. Für $X=\varnothing$ sind die drei Abschlüsse leer.
\end{definition}

\begin{theorem}[Erzeugte Ideale]
Für jede nichtleere Teilmenge $X\subseteq S$ sind
$S^1X$, $XS^1$ und $S^1XS^1$ das kleinste Linksideal,
Rechtsideal beziehungsweise Ideal, das $X$ enthält.
Jedes Links- oder Rechtsideal ist eine Unterhalbgruppe.
\cite{harju,howie}
\end{theorem}

\begin{theorem}[Idealverbände und Idealprodukte]
Die Ideale zusammen mit $\varnothing$ bilden unter Inklusion
einen vollständigen Verband, dessen Infima Durchschnitte und dessen
Suprema Vereinigungen sind. Dasselbe gilt für Links- und Rechtsideale.
Für nichtleere zweiseitige Ideale $I,J$ ist auch $IJ$ ein nichtleeres
Ideal mit $IJ\subseteq I\cap J$.
Insbesondere schneiden sich je zwei nichtleere Ideale.
Jedes unter Inklusion minimale nichtleere Ideal ist bereits
das kleinste nichtleere Ideal und daher eindeutig.
Die Existenz eines solchen Ideals wird nicht behauptet.
\cite{harju,howie}
\end{theorem}

\begin{theorem}[Ideale unter Homomorphismen]
Für einen Homomorphismus $f:S\to T$ ist das Urbild eines Ideals von
$T$ entweder leer oder ein Ideal von $S$.
Das Bild eines Ideals von $S$ ist ein Ideal in der Bildhalbgruppe $f(S)$.
Bei surjektivem $f$ ist es folglich ein Ideal von $T$.
Die entsprechenden Aussagen gelten für Links- und Rechtsideale.
\cite{harju}
\end{theorem}

\section{Zusammenziehen eines Ideals}

\begin{definition}[Rees-Kongruenz und Rees-Quotient]
Für ein Ideal $I$ von $S$ heißt $\rho_I=\Delta_S\cup(I\times I)$ die \emph{Rees-Kongruenz} zu $I$.
Der zugehörige \emph{Rees-Quotient} wird $S/I$ geschrieben.
Er identifiziert alle Elemente von $I$ miteinander und läßt
die übrigen Elemente einzeln bestehen.
\end{definition}

\begin{theorem}[Rees-Quotient und seine Universalität]
\label{b28guide:thm:rees-quotient}
$\rho_I$ ist eine Kongruenz.
Der Quotient $S/I$ kann als $(S\setminus I)\cup\{0_I\}$ dargestellt
werden, wobei $0_I$ ein neues Element ist und
\[
 a\circ b=
 \begin{cases}
 ab,&a,b\in S\setminus I\ \text{und }ab\notin I,\\
 0_I,&\text{sonst}.
 \end{cases}
\]
Ein Homomorphismus $f:S\to T$ faktorisiert genau dann durch $S/I$,
wenn er auf $I$ konstant ist; die Faktorisierung ist eindeutig.
Das neue Nullelement gehört zur Quotientenkonstruktion und muß
nicht von einem Nullelement von $S$ stammen.
\cite{harju,howie}
\end{theorem}

\section{Die allgemeine Hauptfaktorzerlegung}

\begin{definition}[Hauptklasse und Hauptfaktor]
Für $a\in S$ setzen wir
\[
 I(a)=S^1aS^1,\qquad
 J_a=\{b\in S:S^1bS^1=I(a)\},\qquad
 I^-(a)=I(a)\setminus J_a .
\]
$J_a$ ist die in Kapitel~\ref{b28guide:chap:green} betrachtete
$\mathcal J$-Klasse von $a$.
Der \emph{Hauptfaktor} zu $a$ ist $P_a=I(a)$, falls $I^-(a)=\varnothing$, und sonst der Rees-Quotient $P_a=I(a)/I^-(a)$.
Eine nichtleere Halbgruppe heißt \emph{einfach}, wenn sie kein echtes
nichtleeres Ideal besitzt.
Sie heißt \emph{$0$-einfach}, wenn sie ein Nullelement $0$ besitzt,
ihr Quadrat nicht $\{0\}$ ist und ihre einzigen nichtleeren Ideale
$\{0\}$ und die ganze Halbgruppe sind.
Eine \emph{Nullhalbgruppe} besitzt ein Element $0$, auf das jedes
Produkt fällt.
\end{definition}

\begin{theorem}[Hauptfaktor-Trichotomie]
\label{b28guide:thm:principal-factors}
Für jedes $a\in S$ ist $I^-(a)$ entweder leer oder ein Ideal von $S$.
Der Hauptfaktor $P_a$ ist wohldefiniert und gehört zu mindestens
einer der drei Klassen: einfach, $0$-einfach oder Nullhalbgruppe.
Ist $I^-(a)=\varnothing$, so ist $I(a)=J_a$ das kleinste nichtleere
Ideal von $S$ und eine einfache Halbgruppe.
Ist $I^-(a)\ne\varnothing$, so ist $P_a$ entweder $0$-einfach oder
eine Nullhalbgruppe.
Die Trichotomie verlangt keine Endlichkeit und keine Regularität.
Insbesondere wird keine Darstellung aller Hauptfaktoren durch
Rees-Matrizen über Gruppen behauptet.
\cite{howie,bajri}
\end{theorem}
